\section{从谷仓到里巷：西汉晚期的地方世界}
\label{sec:han-late-western-local}

元、成、哀、平之际的长安常以外戚、诏令和受禅的前史出现，仿佛地方社会只在
朝廷失去控制时才忽然进入故事。实际上，谷物、钱币、人户和道路每天都在把中央
的选择翻译成不同地区的成本。关中的仓与市场接近官署，东部王国故地、江淮、岭南
和河西却各有水路、山口、商人、军屯与地方首领。中央能够看到的是郡国送来的
簿籍与奏报；簿籍之外的逃亡、借贷、兼并、婚姻和乡里互助，往往只在冲突发生时
才短暂显影。

元始二年的户口统计因此不应被当作帝国社会的透明横截面。它表明朝廷仍能要求
郡国以户和口的形式申报治理对象，也表明填报需要县吏、里正、土地占有人和迁徙
家庭逐层配合。尹湾文书与居延材料让人看见行政记录的某些环节：收发、核对、
转写和等待。它们不能让我们复原全国每一户的负担，却足以提示一项总数为何会
同时是国家能力的证据和能力边界的证据。一个被迁走的家庭、一个被豪右庇护的
佃客、一个尚未回到原籍的流民，都可能改变簿册与实际生活之间的距离。

\subsection{宫廷的节用与地方的账}

元帝朝裁省宫室、乐府和部分支出的讨论，容易被概括为“节俭”或“衰弱”。对
地方而言，它首先是不同需求在同一财政中竞争：宫廷仪礼、官员俸给、河工、赈贷
和边郡粮饷并不从同一只看不见的口袋中支付。诏令能显示哪些项目被认为可以裁减，
不能说明裁减是否真的抵达每个供役者，也不能把一项中央决定换算为市场价格或
家庭收入。《汉书》帝纪与《食货志》保存的是朝廷的分类；钱币、简牍和地方遗存
保存的是较小尺度的流通与劳作。两种材料不应彼此替代。

北边的呼韩邪关系同样需要付账。赏赐、婚姻与互市不是附在外交辞令后的装饰：
使者要走道路，礼物要入库、出库、转运，边郡要安排接待和防务。昭君故事在后世
获得极强的文学形状，能够确定的是婚姻安排和往来这一政治关系，不能据传说填补
她的意愿、沿途经历或草原各部的内部协商。把这件事放回财政与交通，读者才能
看见“和亲”不是安静的边界，而是一套持续要求中介、马匹、翻译与供给的安排。

\subsection{王氏上升时，谁仍在传递命令}

王氏通过皇后亲属、大司马、尚书与荐举进入中枢，并不意味着其他官署从此停止
运作。相反，正是既有的任命、奏报和近侍通道，使一个亲属网络能够在不必公开
取代所有制度的情况下变得更有分量。等待迁调的官员、需要核验赋役的地方吏员、
希望获得赈贷或军事支援的郡守，都必须判断消息经谁递送才会被看见。外戚政治
因此不是宫门内的附属剧情；它改变了地方请求怎样抵达皇帝，也改变了谁能把中央
资源变成私人声望。

这种集中并非只产生压迫，也常提供短期的协调能力。幼主、灾荒或边警发生时，
能够迅速调动文书和近臣的集团可以让命令不停摆。代价在于，其他行动者愈难确认
自己的申诉会被同样对待。王莽居摄前的符命与劝进，正是在这一不均衡中获得力量：
它们把中枢已有的控制写成普遍同意，却无法保证市场、乡里、边郡或正在迁徙的
人群愿意照新名称重新安排生活。

\begin{sourcenote}[晚西汉地方史能说到哪里]
《汉书》帝纪、外戚传、地理志与《食货志》可追踪朝廷的户口、财政和任命语言；
尹湾、居延等文书以及钱币、墓葬和区域考古提供局部物质尺度。它们都没有留下
连续的全国性家庭账簿。本节因此只说明簿籍、财政和中介如何构成压力，不把统计
数字、后世昭君叙事或外戚道德评价改写成地方社会的一致经验。
\end{sourcenote}
